% Figure: stress element for the combined bending+torsion worked example
% on a shaft. Shows a surface element with a normal (bending) stress on
% the x-faces and a shear (torsion) stress on all faces, then the Mohr
% circle reduction to principal stresses for sigma_x=120, tau=50 MPa.
\begin{figure}[htbp]
\centering
\begin{tikzpicture}[scale=1.0]
% === left: the stress element ===
\begin{scope}[shift={(-4.2,0)}]
\draw[engblack, very thick] (-1,-1) rectangle (1,1);
% bending normal stress on x-faces (horizontal arrows)
\draw[engblue, very thick, -{Latex}] (1,0) -- (2,0);
\draw[engblue, very thick, -{Latex}] (-1,0) -- (-2,0);
\node[engblue, font=\scriptsize, anchor=west] at (2,0) {$\sigma_x = 120$};
% torsion shear on the faces (tangential arrows)
\draw[engred, very thick, -{Latex}] (-1,1) -- (0,1);
\draw[engred, very thick, -{Latex}] (1,-1) -- (0,-1);
\draw[engred, very thick, -{Latex}] (1,1) -- (1,0);
\draw[engred, very thick, -{Latex}] (-1,-1) -- (-1,0);
\node[engred, font=\scriptsize, anchor=south] at (-0.5,1.05) {$\tau = 50$};
\node[engblack, font=\scriptsize, anchor=north] at (0,-1.6) {surface element (MPa)};
\end{scope}
% === right: the Mohr circle ===
\begin{scope}[shift={(2.6,0)}]
\begin{axis}[
  width=7.5cm, height=6.5cm,
  axis equal image,
  xlabel={$\sigma$ (MPa)},
  ylabel={$\tau$ (MPa)},
  xmin=-40, xmax=160, ymin=-90, ymax=90,
  axis lines=middle,
  tick label style={font=\tiny},
  label style={font=\scriptsize},
]
% center C = (sigma_x/2, 0) = (60,0); R = sqrt(60^2 + 50^2) = 78.1
\draw[engblue, very thick] (axis cs:60,0) circle[radius=78.1];
\filldraw[engred] (axis cs:120,50) circle[radius=1.6pt];
\node[engred, font=\tiny, anchor=south west] at (axis cs:120,50) {$X(120,50)$};
\filldraw[engred] (axis cs:0,-50) circle[radius=1.6pt];
\node[engred, font=\tiny, anchor=north east] at (axis cs:0,-50) {$Y(0,-50)$};
\draw[engred, thick] (axis cs:120,50) -- (axis cs:0,-50);
\filldraw[engorange] (axis cs:138.1,0) circle[radius=1.6pt];
\node[engorange, font=\tiny, anchor=south] at (axis cs:138,6) {$\sigma_1=138$};
\filldraw[engorange] (axis cs:-18.1,0) circle[radius=1.6pt];
\node[engorange, font=\tiny, anchor=south] at (axis cs:-22,6) {$\sigma_2=-18$};
\end{axis}
\end{scope}
\end{tikzpicture}
\caption[Combined bending and torsion on a shaft surface]{The surface
element of a shaft under combined bending and torsion (left) carries a
normal stress $\sigma_x = 120\ \text{MPa}$ from bending and a shear
$\tau = 50\ \text{MPa}$ from torque. The Mohr circle (right) has centre
$C = (\sigma_x/2, 0) = (60, 0)$ and radius $R = \sqrt{60^2 + 50^2} =
78.1\ \text{MPa}$, giving principal stresses $\sigma_1 = 138$ and
$\sigma_2 = -18\ \text{MPa}$ and a maximum in-plane shear $\tau_{\max} =
78\ \text{MPa}$. The compressive minor principal stress raises the
Tresca and von Mises equivalents above the bending stress alone.}
\label{fig:vol03ch03:combined-loading-element}
\end{figure}
